\documentclass[a4paper,12pt]{article}

\usepackage[utf8]{inputenc}
\usepackage[T1]{fontenc}
\usepackage[italian]{babel}
\usepackage{geometry}
\usepackage{hyperref}

\geometry{margin=2.5cm}

\title{Differenze tra versione locale e versione LAN}
\author{AFAM Connect}
\date{}

\begin{document}

\maketitle

\section{Differenze tra versione locale e versione LAN}

La versione originale del progetto, contenuta nella cartella \texttt{/IDS/Codice},
rappresenta l'applicazione desktop locale classica. Ogni installazione utilizza il
proprio database SQLite locale, indicato dal file \texttt{afam.db}, e i contenuti
multimediali vengono riferiti tramite percorsi assoluti del computer su cui sono
stati caricati.

La versione sperimentale contenuta nella cartella \texttt{Codice-LAN-Test}, invece,
mantiene l'architettura basata sul pattern Repository, ma introduce una
configurazione condivisa per consentire a piu postazioni presenti nella stessa LAN
di accedere agli stessi dati e agli stessi file.

\subsection{Architettura}

L'architettura rimane basata sul Repository: i sottosistemi applicativi continuano
a comunicare con il database esclusivamente attraverso il componente
\texttt{DatabaseManager}.

\begin{center}
\texttt{Interfaccia JavaFX -> Control -> DatabaseManager -> SQLite}
\end{center}

Nella versione LAN non viene introdotto un server applicativo. Le diverse
postazioni eseguono la stessa applicazione e puntano allo stesso database
condiviso all'interno della rete locale del nodo AFAM.

\begin{center}
\texttt{Postazione 1 -> DatabaseManager -> DB condiviso}

\texttt{Postazione 2 -> DatabaseManager -> DB condiviso}
\end{center}

Pertanto la versione sperimentale puo essere descritta come una architettura
Repository con database condiviso su LAN, non come una architettura client-server
applicativa.

\subsection{Database SQLite condiviso}

Nella versione originale il componente \texttt{BoundaryDBMS} apre sempre il
database locale:

\begin{verbatim}
jdbc:sqlite:afam.db
\end{verbatim}

Nella versione LAN il percorso del database e configurabile tramite proprieta di
sistema o variabile ambiente:

\begin{verbatim}
-Dafam.db.path=PERCORSO_DATABASE
AFAM_DB_PATH=PERCORSO_DATABASE
\end{verbatim}

Se nessun percorso viene specificato, l'applicazione continua a usare il database
locale \texttt{afam.db}. In questo modo la modifica e retrocompatibile con il
funzionamento originale.

\subsection{Cartella condivisa per i file}

Nella versione originale, foto profilo, immagini, audio, video e documenti non
vengono salvati direttamente nel database. Nel database viene memorizzato solo il
percorso fisico del file, ad esempio:

\begin{verbatim}
C:\Users\Simone\Desktop\video.mp4
\end{verbatim}

Questo approccio funziona su un singolo computer, ma non consente a un altro
computer della LAN, ad esempio un Mac, di aprire lo stesso file.

Per questo nella versione LAN e stato introdotto il componente
\texttt{ArchivioFile}. Tale componente copia i nuovi file caricati in una cartella
condivisa e salva nel database un percorso relativo.

La struttura consigliata e:

\begin{verbatim}
AFAM-LAN
|-- DB
|   |-- afam.db
|-- FILES
    |-- foto
    |-- contenuti
\end{verbatim}

Su Windows l'applicazione viene avviata con:

\begin{verbatim}
-Dafam.db.path=C:\Users\Simone\Desktop\AFAM-LAN\DB\afam.db
-Dafam.files.path=C:\Users\Simone\Desktop\AFAM-LAN\FILES
\end{verbatim}

Su Mac, dopo aver montato la condivisione SMB, viene avviata con:

\begin{verbatim}
-Dafam.db.path=/Volumes/AFAM-LAN/DB/afam.db
-Dafam.files.path=/Volumes/AFAM-LAN/FILES
\end{verbatim}

Nel database non viene piu salvato un percorso assoluto del singolo computer, ma
un percorso relativo, ad esempio:

\begin{verbatim}
contenuti/video_uuid.mp4
foto/profilo_uuid.png
\end{verbatim}

Ogni postazione ricostruisce poi il percorso reale usando il proprio valore di
\texttt{afam.files.path}.

\subsection{File modificati}

Rispetto alla versione originale, nella versione LAN sono stati modificati o
aggiunti i seguenti file principali:

\begin{itemize}
    \item \texttt{pom.xml}: aggiunte le proprieta \texttt{afam.db.path} e
    \texttt{afam.files.path}.
    \item \texttt{BoundaryDBMS.java}: reso configurabile il percorso del database
    SQLite.
    \item \texttt{ArchivioFile.java}: nuovo componente per archiviare e risolvere
    i file condivisi.
    \item \texttt{UploadFileControl.java}: copia i nuovi contenuti nella cartella
    condivisa.
    \item \texttt{GestioneFileControl.java}: gestisce la modifica dei contenuti
    usando la cartella condivisa.
    \item \texttt{GestioneProfiloControl.java}: archivia le foto profilo nella
    cartella condivisa.
    \item \texttt{SchermataDatiPersonali.java}: visualizza la foto profilo
    risolvendo il percorso condiviso.
    \item \texttt{SchermataProfiloStudente.java}: visualizza la foto dello
    studente consultato tramite percorso condiviso.
    \item \texttt{SchermataVisualizzatoreDocumenti.java}: apre immagini e PDF
    usando il percorso risolto.
    \item \texttt{SchermataRiproduttoreMultimediale.java}: apre audio e video
    usando il percorso risolto.
    \item \texttt{DatabaseManager.java}: calcola file fisici e dimensioni tramite
    \texttt{ArchivioFile}.
    \item \texttt{README.md}: documenta la configurazione LAN Windows/Mac.
\end{itemize}

\subsection{Comportamento ottenuto}

La versione LAN consente a due computer nella stessa rete locale di:

\begin{itemize}
    \item accedere allo stesso database SQLite;
    \item visualizzare profili e background artistico degli altri utenti;
    \item visualizzare contenuti pubblici caricati da un'altra postazione;
    \item condividere foto, immagini, audio, video e documenti tramite una
    cartella comune;
    \item mantenere il pattern Repository come punto unico di accesso ai dati.
\end{itemize}

I contenuti privati restano non visibili nella sezione pubblica fino a quando non
vengono impostati come pubblici oppure condivisi tramite codice di sblocco.

\subsection{Limiti della simulazione}

La soluzione mantiene SQLite e non introduce un server applicativo. Questo la rende
coerente con l'architettura Repository gia documentata, ma presenta alcuni limiti:

\begin{itemize}
    \item SQLite su cartella condivisa e adatto a una simulazione LAN o a bassa
    concorrenza;
    \item scritture contemporanee da piu postazioni potrebbero causare lock;
    \item per un uso produttivo sarebbe preferibile un DBMS server;
    \item i file caricati prima della modifica mantengono vecchi percorsi assoluti
    e devono essere ricaricati o migrati nella cartella condivisa.
\end{itemize}

In conclusione, la versione \texttt{Codice-LAN-Test} estende il progetto originale
per una simulazione LAN mantenendo l'architettura Repository: non viene introdotta
una comunicazione diretta tra applicazioni, ma piu istanze dell'applicazione
condividono lo stesso database e la stessa cartella file.

\end{document}
